% SOURCE_AUTHORITY: sga5_fr_workpass.tex, lines 11729--11769 (LF-normalized unit).
% SOURCE_SHA256: 791F4EFFC5E02832D5D77ED03518C8156D6F07E4C8238B03545DB93D883FBB28
\subsection*{3. Carácter funtorial}

Sean $\cC$ y $\cC'$ dos categorías trianguladas, y sea $T:\cC\to\cC'$ un funtor exacto de $\cC$ en $\cC'$, es decir, un funtor aditivo, «graduado» y que transforma los triángulos distinguidos en triángulos distinguidos. Denotemos, como antes, por
\[
A,A':\Ab\longrightarrow\Ab
\]
los funtores asociados a las categorías $\cC$ y $\cC'$. Resulta un morfismo funtorial $A'\to A$ de «composición con $T$» y, por tanto, un homomorfismo
\[
K(T):K(\cC)\longrightarrow K(\cC')
\]
único tal que el diagrama
\[
\begin{tikzcd}[column sep=large,row sep=large]
\Ob\cC \arrow[r,"T"] \arrow[d,"\operatorname{cl}_{\cC}"']&
\Ob\cC' \arrow[d,"\operatorname{cl}_{\cC'}"]\\
K(\cC) \arrow[r,"K(T)"']& K(\cC')
\end{tikzcd}
\]
sea conmutativo. Resulta que, si los funtores $T,T_1:\cC\to\cC'$ son isomorfos, entonces $K(T)=K(T_1)$. En particular, si el funtor $T$ es una equivalencia de categorías, entonces $K(T)$ es un isomorfismo.

Si $\cC''$ es una tercera categoría triangulada y
\[
F:\cC\times\cC'\longrightarrow \cC''
\]
un bifuntor exacto en cada variable, entonces es claro que resulta una aplicación bilineal
\[
K(F):K(\cC)\times K(\cC')\longrightarrow K(\cC''),
\]
tal que se verifica la relación siguiente:
\[
K(F)(\operatorname{cl}_{\cC}(X),\operatorname{cl}_{\cC'}(X'))
=
\operatorname{cl}_{\cC''}(F(X,X')).
\]

\begin{statement}{Proposición 3.1}
Sea $\cA$ una categoría triangulada, $\cB$ una subcategoría triangulada gruesa de $\cA$, $\cA/\cB$ la categoría cociente de $\cA$ por $\cB$, $Q:\cA\to\cA/\cB$ el funtor canónico de paso al cociente, e $i:\cB\to\cA$ el funtor de inclusión. Entonces la sucesión siguiente es exacta:
\[
K(\cB)\xrightarrow{K(i)}K(\cA)\xrightarrow{K(Q)}K(\cA/\cB)\longrightarrow 0.
\]
\end{statement}
